% Transcription — fonds Alexandre Grothendieck, shelfmark 1, batch 7 (pages 121-140).
% Established from the facsimile archives/batches/1/batch-07.pdf (a mirror of
% https://grothendieck.umontpellier.fr/1.pdf), per the skill
% transcribe-grothendieck.
%
% Pass: Fable 5.1 (claude-fable-5-1), 2026-09-10 — first pass, unchecked against
% the pages by a human.
%
% Eighteen of the twenty pages are transcribed; pages 122 and 140 are blank.
% \page{} numbers the position in the folder, as batches 1 to 6 did: the
% archivists' pencil numbers run two ahead, so page 121 here is pencilled
% 123 and page 139 is pencilled 141.
%
% The batch falls in three parts.
%
%   Page 121 — the end of « Sur une certaine fonction holomorphe sur l'espace
%     de Banach ℓ¹ », continuing the proof broken off at the foot of page 120
%     (batch 6): ⁿ√(n! α_n) → 0, the n-linear form α_n(λ^(1),…,λ^(n)) =
%     (1/n!) Σ λ^(1)_{i_1}⋯λ^(n)_{i_n} over distinct indices, ‖α_n‖ ≤ 1/n!, the
%     radius R_0 ≥ 1, and the elements a^(m) = (1/m,…,1/m,0,…) with
%     n! α_n(a^(m)) → 1 showing S is unbounded on the unit ball, whence
%     R_0 = R_1 = 1.
%   Pages 123-125 — older sheets in a slower hand: the end of an inequality
%     (‖u_1∧⋯∧u_n‖_1)² ≤ α_n(|u_1|,…,|u_n|) α_n(|u_1|',…,|u_n|') proved from
%     the canonical representations u_i = Σ λ a ⊗ b and Cauchy-Schwarz, then
%     five boxed corollaries: |α_n(u_1,…,u_n)| ≤ ‖u_1∧⋯∧u_n‖_1 ≤ ½(…+…),
%     |α_n(A+B)|² ≤ α_n(|A|+|B|) α_n(|A|'+|B|'), |det(1+A+B)| ≤ ½(det(1+|A|+|B|)
%     + det(1+|A|'+|B|')), |det(1+A+B)| ≤ det(1+|A|) det(1+|B|) reduced to the
%     hermitian positive case through det(1+HK) ≤ det(1+K) for ‖H‖ ≤ 1 and
%     Λⁿ(HK) = (ΛⁿH)(ΛⁿK), and det(1+|A+B|) ≤ det(1+|A|) det(1+|B|). Page 125
%     is the same leaf as 124 with a slip pasted over its lower half; only
%     the slip — the derivation of corollary 5 from corollary 4 through
%     |A+B| = U(A+B) — is transcribed under \page{125}.
%   Pages 126-139 — « Théorème » on the classes 𝓛^p(H), 0 < p ≤ ∞, and its
%     proof, numbered (1) to (14). Page 126 states it: 𝓛^p(H) a two-sided
%     ideal stable under A → A*, (1) ‖A‖_p = ‖A*‖_p, (2) ‖BA‖_p ≤ ‖B‖‖A‖_p,
%     (3) the triangle inequality for p ≥ 1 with completeness, (4)
%     S_p(A+B) ≤ S_p(A)+S_p(B) for S_p = ‖·‖_p^p and 0 < p ≤ 1, giving a
%     complete metrisable non-locally-convex space, (5) Σ|λ_i|^p ≤ S_p(A) for
%     the eigenvalues, monotonicity on positive operators, and (6)
%     ‖AB‖_r ≤ ‖A‖_p‖B‖_q for 1/p+1/q ≥ 1/r. Pages 127-128: (1) and (2)
%     from the singular values, (3) from (6) through (7) ‖A‖_p =
%     Sup_{‖B‖_{p'} ≤ 1}|Tr AB|, and Lemma 1, (8) Σ α_i^p = (p sin πp/π)
%     ∫ log f(v)/v^{1+p} dv for f = ∏(1+α_i v). Pages 129-131: the corollary
%     (9) S_p(A) = (p sin πp/π) ∫ log M_{|A|}(v)/v^{1+p} dv with M_A(v) =
%     Sup_{|z|=v} det(1+zA), (4) from (10) M_{|A+B|} ≤ M_{|A|}M_{|B|}, and (5)
%     from the classical bound s_p(A) ≤ p² ∫ log M_A/r^{p+1} for the zeros
%     of det(1−zA), giving s_p(A) ≤ (pπ/sin pπ) S_p(A), then s_p(A) =
%     s_{p/n}(Aⁿ) and n → ∞; the lower half of page 131, a first attack on
%     (6), is cancelled by diagonal strokes. Pages 132-135: (11)
%     ‖A_1⋯A_{2ⁿ}‖_1 ≤ ∏‖A_i‖_{2ⁿ} by induction on n through C =
%     (ABB*A*)^{2^{n−2}} and Tr C = Tr C', Hölder ‖AB‖_1 ≤ ‖A‖_p‖B‖_q via the
%     polar decompositions and dyadic exponents λ = 2ⁿ/α, and the lemma (12)
%     α_n(L^m) ≤ α_n(H^m) α_n(K^m) for H = A*A, K = B*B, L = (AB)*(AB).
%     Pages 136-139: the general case of (6) with 1/p+1/q = 1/r, m with
%     p, q ≤ 2m, S_r(AB) = S_{r/2m}(L^m) written through (9), M_{L^m}(ρ) ≤
%     M_{H^m}(ρ^{r/p}) M_{K^m}(ρ^{r/q}), the substitution σ = ρ^{r/p} giving
%     (13) λ(r/2m) S_r(AB) ≤ λ(p/2m) S_p(A) + λ(q/2m) S_q(B) with λ(x) =
%     π/sin πx, the limit (14) S_r(AB) ≤ (r/p) S_p(A) + (r/q) S_q(B), and the
%     homogeneity argument A → λA, B → μB with the Lagrange minimisation of
%     αax + βby on x^α y^β = 1, ending at c ≤ a^α b^β on page 139.
%
% Cancellations: nearly every page from 127 on carries struck words or lines,
% transcribed as \struck{} with one \note{} per block where the block cannot
% be read; pages 131 and 135 carry whole cancelled passages. Boxed formulas
% are noted as « encadré ». Notes in his hand in the margins of pages 121,
% 126, 127, 132, 134, 136 and 138 are transcribed as \marginal{} where they
% can be read and noted where they cannot. Page 126 opens with four pencil
% lines, partly struck, of which only a few words are read.
%
% Notation: 𝓛^p(H) set \mathcal{L}^p(H) (page 126 twice writes 𝓛^p(A),
% kept); ‖·‖_p the Schatten norm and S_p(A) = ‖A‖_p^p; s_p(A) = Σ|λ_i|^p
% over the eigenvalues; α_n(A) the coefficients of det(1+zA) and M_A(v) =
% Sup_{|z|=v} det(1+zA); |A| = (A*A)^{1/2} and |A|' its counterpart for A*,
% as in batches 4 to 6; Λⁿ the exterior power, whose exponent he writes
% above the letter; 2ⁿ written « 2n » on pages 132-133, set 2^n on the
% strength of the recurrence; Sup, Inf, Tr set as he writes them; « cpt »
% for compact and « top. » for topologique kept; « Cauchy-Schwartz » as he
% spells it. Formula numbers are his and restart with the theorem on page
% 126; (10) and (11) are each used twice.
%
% The dating is the inventory's own for the shelfmark. Its group,
% « MATHÉMATIQUES / RECHERCHE / Travaux », carries no date.
\documentclass[11pt,a4paper]{article}
% Le préambule commun à toutes les transcriptions du fonds Grothendieck.
%
% Il tient en une idée : **l'incertitude se compose, elle ne se résout pas.**
% Une transcription qui lisse un mot douteux a détruit la seule chose pour
% laquelle on la faisait — un lecteur doit pouvoir savoir, à chaque endroit,
% ce qui était sur le feuillet et ce qui a été ajouté. Les six macros qui
% suivent sont donc l'appareil critique, et non de la décoration.
%
% Les mêmes commandes sont reconnues par `scripts/render.mjs`, qui produit la
% vue de lecture du site. Une macro ajoutée ici doit l'être aussi là-bas,
% faute de quoi le rendu échouera bruyamment — ce qui est voulu : un rendu
% silencieusement faux est pire qu'un rendu absent.

\ProvidesPackage{grothendieck}[2026/08/08 Transcriptions du fonds Grothendieck]

\usepackage{amsmath,amssymb,amsthm}
\usepackage{mathrsfs}
\usepackage{stmaryrd}
% Les diagrammes commutatifs sont partout dans ce fonds : c'est la langue
% même de la géométrie algébrique de Grothendieck, et une transcription qui
% les rendrait en prose ne transcrirait rien.
\usepackage{tikz-cd}
% Français par défaut — c'est la langue des feuillets — mais l'anglais est
% chargé aussi : la traduction et le résumé sont des documents anglais, et la
% typographie française (espaces avant les deux-points) y serait une faute.
% Ces documents basculent par \selectlanguage{english} en tête de corps.
\usepackage[english,french]{babel}
\usepackage{fontspec}
\usepackage{geometry}
\geometry{a4paper,margin=25mm,marginparwidth=28mm}
\usepackage{marginnote}
\usepackage{xcolor}
% ulem, not soul. `soul`'s \st and \ul are built on a character-by-character
% scan that XeLaTeX's font machinery breaks: the document fails with an
% undefined control sequence rather than degrading. ulem does the same two jobs
% and is Unicode-engine safe. `normalem` stops it hijacking \emph, which the
% transcriptions use for Grothendieck's own emphasis.
\usepackage[normalem]{ulem}
\usepackage{hyperref}
\hypersetup{colorlinks=true,linkcolor=black,urlcolor=black,pdfborder={0 0 0}}

% --- Métadonnées du lot -----------------------------------------------------
% Reprises telles quelles par le script de rendu, qui les affiche en tête de
% la vue de lecture. Elles fixent ce que le fichier prétend couvrir : sans
% elles, une transcription flotte sans point d'attache dans un fonds de seize
% mille feuillets.

\newcommand{\folder}[1]{\gdef\grothFolder{#1}}
\newcommand{\batch}[1]{\gdef\grothBatch{#1}}
\newcommand{\leaves}[2]{\gdef\grothFirst{#1}\gdef\grothLast{#2}}
\newcommand{\foldertitle}[1]{\gdef\grothFoldertitle{#1}}
\gdef\grothFolder{?}\gdef\grothBatch{?}\gdef\grothFirst{?}\gdef\grothLast{?}\gdef\grothFoldertitle{}

% --- Le résumé --------------------------------------------------------------
%
% Il ouvre la lecture modernisée, sous le titre « Résumé », et il est composé
% en petit. Ce n'est pas un ornement : c'est le seul passage du document qui
% ne soit pas des mathématiques, et un lecteur doit pouvoir le reconnaître
% avant de l'avoir lu — puis le sauter s'il connaît déjà le sujet, ou s'y
% arrêter s'il ne le connaît pas. Le filet qui le suit marque la reprise du
% corps.

\newenvironment{resume}
  {\small\setlength{\parskip}{0.45em}}
  {\par\medskip\hrule\medskip}

% --- Mots-clés ---------------------------------------------------------------
%
% En anglais, en clôture du résumé de la lecture modernisée : le vocabulaire
% moderne sous lequel chercher ce que les feuillets construisent. Le manifeste
% du site (`npm run manifest`) les extrait pour étiqueter la cote entière —
% cette ligne est la source unique des tags, il n'y a pas de fichier à côté.
\newcommand{\keywords}[1]{\par\smallskip\noindent
  {\footnotesize\textsc{Keywords} --- \textit{#1}}\par}

% --- L'appareil critique ----------------------------------------------------

% Le numéro de feuillet, en marge. C'est l'ancre qui permet de régler un
% désaccord sur une formule en regardant *un* feuillet plutôt qu'une chemise
% de six cents. Le site s'en sert aussi pour tourner les pages du fac-similé
% quand on fait défiler la transcription.
\newcommand{\leaf}[1]{%
  \marginnote{\footnotesize\color{gray!70}\textsf{#1}}%
  \label{leaf:#1}\ignorespaces}

% Illisible : on ne devine pas. Un mot inventé qui se lit comme les autres est
% le pire résultat possible.
\newcommand{\ill}{\textcolor{red!60!black}{[\,\ldots\,]}}

% Lecture douteuse : le mot est donné, et signalé comme incertain.
\newcommand{\uncertain}[1]{\uline{#1}}

% Ajout de l'éditeur — une lettre manquante, un mot sous-entendu.
\newcommand{\add}[1]{\textcolor{blue!50!black}{[#1]}}

% Biffé par Grothendieck lui-même. Ce qu'il a rayé fait partie du document :
% une rature dit souvent où la pensée a changé de direction.
\newcommand{\struck}[1]{\sout{#1}}

% Note du transcripteur, sur l'état matériel du feuillet.
\newcommand{\note}[1]{\par\smallskip{\small\color{gray!60!black}\textit{[#1]}}\par\smallskip}

% Note marginale de Grothendieck, distincte du corps du texte.
\newcommand{\marginal}[1]{\par\smallskip\noindent\hspace{1em}%
  {\small\color{gray!60!black}#1}\par\smallskip}

% --- Titre ------------------------------------------------------------------

\AtBeginDocument{%
  \begin{center}
    {\large\bfseries Fonds Alexandre Grothendieck --- cote n\textsuperscript{o}~\grothFolder}\\[2pt]
    {\itshape\grothFoldertitle}\\[4pt]
    {\small feuillets \grothFirst--\grothLast\ (lot \grothBatch)}\\[2pt]
    {\footnotesize\color{gray!60!black}Transcription établie d'après le fac-similé numérisé
     par l'Université de Montpellier.\\
     Les lectures incertaines sont soulignées, les illisibles notées [\,\ldots\,],
     les ajouts entre crochets.}
  \end{center}
  \vspace{1em}}

\endinput


\folder{1}
\batch{7}
\pages{121}{140}
\dating{1953}
\watermark{Édition de démonstration}
\foldertitle{Analyse fonctionnelle : notes manuscrites (s.d.), lettre (1953)}

\begin{document}

\note{La page 121 achève la démonstration commencée au bas de la page 120 du lot précédent. Les pages 123 à 125 sont des feuillets plus anciens, d'une main plus lente, sur les puissances extérieures et le déterminant de Fredholm. Les pages 126 à 139 sont d'un seul tenant : l'énoncé d'un théorème sur les classes $\mathcal{L}^p(H)$ et sa démonstration, dont les formules sont numérotées de (1) à (14).}

\page{121}
\note{la page poursuit la démonstration de la proposition commencée au bas de la page 120}
\[
  n!\,\alpha_n \leq \frac{\varepsilon_0}{n+1}\,\varepsilon_1 \cdots \varepsilon_n\, e^{n}
  \qquad
  \sqrt[n]{n!\,\alpha_n} \leq e\Bigl(\frac{\varepsilon_0}{n+1}\Bigr)^{\frac1n} (\varepsilon_1 \cdots \varepsilon_n)^{\frac1n}
\]
\note{comme à la page 120, les indices et les exposants se lisent $n$ ou $n+1$ ; le dernier facteur est réécrit au-dessus d'un facteur biffé} $(n \to \infty)$, $\bigl(\frac{\varepsilon_0}{n+1}\bigr)^{\frac1n} \to 1$, et $(\varepsilon_1 \cdots \varepsilon_n)^{\frac1n} \to 0$, \uncertain{d'où} \ill{} $\sum n!\,\alpha_n(\lambda)$ \uncertain{converge} \uncertain{absolument} \struck{\ill{}}, \uncertain{d'où} $\sqrt[n]{n!\,\alpha_n} \to 0$. \ill{}

\marginal{\ill{} $S(\lambda)$ \ill{}} \note{note marginale en biais, dans le coin supérieur gauche, reliée par une accolade aux deux lignes qui suivent ; illisible hormis $S(\lambda)$}

Comme on a $|\alpha_n(\lambda)| \leq \|\lambda\|_1^{\,n}$, il est \ill{} \uncertain{que} la fonction $S(\lambda)$ est holomorphe au voisinage de l'origine ; d'ailleurs, \uncertain{la} \uncertain{forme} \ill{} \uncertain{linéaire} $\alpha_n(\lambda^{(1)}, \ldots, \lambda^{(n)})$ \uncertain{associée} à $\lambda \mapsto \alpha_n(\lambda)$ \uncertain{est}
\[
  \alpha_n(\lambda^{(1)}, \ldots, \lambda^{(n)}) = \frac{1}{n!} \sum_{i_1, \ldots, i_n} \lambda^{(1)}_{i_1} \cdots \lambda^{(n)}_{i_n}
\]
(la somme étant étendue aux systèmes $(i_1, \ldots, i_n)$ d'indices \uncertain{tous} distincts) d'où $\|\alpha_n\| \leq \dfrac{1}{n!}$ \struck{$\|n!\,\alpha_n\| \leq 1$}, \uncertain{la} \uncertain{marge} \uncertain{de} \ill{}. \uncertain{Le} \uncertain{rayon} $R_0$ de \uncertain{la} \uncertain{fonction} $S(\lambda)$ à l'origine est donc $\geq 1$. Soit alors $a^{(m)}$ l'élément de $\ell^1$ \uncertain{défini} \uncertain{par}
\[
  a^{(m)}_i = \begin{cases} \dfrac{1}{m} & \text{si } i \leq m \\[4pt] 0 & \text{si } i > m \end{cases}
\]
\uncertain{on} \uncertain{a} \uncertain{manifestement} $\|a^{(m)}\|_1 = 1$, et
\[
  \alpha_n(a^{(m)}) = \struck{\ill{}}\ \frac{1}{m^n}\binom{m}{n} = \frac{1}{m^n}\,\frac{m!}{n!\,(m-n)!} \quad \text{si } n \leq m,
\]
\note{le coefficient binomial est écrit $C_m^n$ ; un premier membre de droite est biffé} d'où
\[
  n!\,\alpha_n(a^{(m)}) = \frac{1}{m^n}\, m(m-1)\cdots(m-n+1) = \Bigl(1 - \frac{1}{m}\Bigr)\Bigl(1 - \frac{2}{m}\Bigr)\cdots\Bigl(1 - \frac{n-1}{m}\Bigr)
\]
Donc $n!\,\alpha_n(a^{(m)}) \to 1$ pour $m \to \infty$. Il suit que $S(a^{(m)}) \to \infty$ si $m \to \infty$ ; donc $S(\lambda)$ \uncertain{n'est} \uncertain{pas} \uncertain{bornée} dans la boule unité de $\ell^1$, et \uncertain{par} \uncertain{conséquent} \uncertain{le} \uncertain{rayon} \ill{} \uncertain{de} \uncertain{convergence} \ill{} $\leq 1$, d'où $R_0 = R_1 = 1$. \note{« $R_1$ » est repris au-dessus de la ligne, avec « bref »} Enfin, $S(\lambda)$ \uncertain{étant} \uncertain{holomorphe} \uncertain{pour} \ill{} \uncertain{et} \uncertain{prenant} \ill{} \uncertain{holomorphes} \uncertain{sur} \uncertain{chaque} \ill{} \note{les deux dernières lignes sont encadrées d'un trait ; la fin de la phrase est perdue}

\page{123}
\note{la page s'ouvre au milieu d'un énoncé dont le début manque ; feuillets plus anciens, d'une main plus lente}
\[
  \bigl(\|u_1 \wedge \cdots \wedge u_n\|_1\bigr)^2 \leq \alpha_n(|u_1|, \ldots, |u_n|)\; \alpha_n(|u_1|', \ldots, |u_n|')
\]
\[
  |\alpha_n(u_1, \ldots, u_n)| \leq \alpha_n(|u_1|, \ldots, |u_n|).
\]
\note{le second facteur de la première ligne est réécrit au-dessus d'une écriture biffée}

\emph{Démonstration} Posons $u_i = \sum_\alpha \lambda^{(i)}_\alpha\, a^{(i)}_\alpha \otimes b^{(i)}_\alpha$ (représentation canonique de $u_i$) alors
\[
  |u_i| = \sum_\alpha \lambda^{(i)}_\alpha\, a^{(i)}_\alpha, \qquad |u_i|' = \sum_\alpha \lambda^{(i)}_\alpha\, b^{(i)}_\alpha
\]
\note{ainsi sur la page, sans les produits tensoriels $a \otimes a$, $b \otimes b$}
\[
  u_1 \wedge \cdots \wedge u_n = \frac{1}{n!} \sum_{\alpha_1 \ldots \alpha_n} \lambda^{(1)}_{\alpha_1} \cdots \lambda^{(n)}_{\alpha_n}\, (a^{(1)}_{\alpha_1} \wedge \cdots \wedge a^{(n)}_{\alpha_n}) \otimes (b^{(1)}_{\alpha_1} \wedge \cdots \wedge b^{(n)}_{\alpha_n})
\]
d'où
\[
  \|u_1 \wedge \cdots \wedge u_n\|_1 \leq \frac{1}{n!} \sum_{\alpha_1 \ldots \alpha_n} \lambda^{(1)}_{\alpha_1} \cdots \lambda^{(n)}_{\alpha_n}\, \|a^{(1)}_{\alpha_1} \wedge \cdots \wedge a^{(n)}_{\alpha_n}\|\, \|b^{(1)}_{\alpha_1} \wedge \cdots \wedge b^{(n)}_{\alpha_n}\|
\]
\uncertain{d'autre} \uncertain{part} \ill{} \note{une ligne courte, interlinéée de « $= \operatorname{Tr}(u_1 \wedge \cdots \wedge u_n)$ » lu avec doute, et se terminant par un $2$ ou $\ell$ suivi de « $= \alpha$ … »}
\[
  \alpha_n(|u_1|, \ldots, |u_n|) = \frac{1}{n!} \sum_{\alpha_1 \ldots \alpha_n} \lambda^{(1)}_{\alpha_1} \cdots \lambda^{(n)}_{\alpha_n}\, \|a^{(1)}_{\alpha_1} \wedge \cdots \wedge a^{(n)}_{\alpha_n}\|^2 \struck{\ill{}}
\]
\[
  \alpha_n(|u_1|', \ldots, |u_n|') = \frac{1}{n!} \sum_{\alpha_1 \ldots \alpha_n} \lambda^{(1)}_{\alpha_1} \cdots \lambda^{(n)}_{\alpha_n}\, \|b^{(1)}_{\alpha_1} \wedge \cdots \wedge b^{(n)}_{\alpha_n}\|^2
\]
L'inégalité \uncertain{cherchée} \uncertain{résulte} alors \uncertain{aussitôt} de l'inégalité \struck{\ill{}} de Cauchy-Schwartz. \note{une ligne entière est biffée avant « de Cauchy-Schwartz » ; « Schwartz » ainsi orthographié}

\emph{Corollaire 1} \note{encadré}
\[
  |\alpha_n(u_1, \ldots, u_n)| \leq \|u_1 \wedge \cdots \wedge u_n\|_1 \leq \tfrac12\bigl(\alpha_n(|u_1|, \ldots, |u_n|) + \alpha_n(|u_1|', \ldots, |u_n|')\bigr)
\]

\emph{Corollaire 2} \note{encadré}
\[
  |\alpha_n(A+B)|^2 \leq \alpha_n(|A| + |B|)\; \alpha_n(|A|' + |B|')
\]
En effet, on a \struck{\ill{}}
\begin{gather*}
  |\alpha_n(A+B)|^2 = \Bigl|\sum_k C_n^k\, \alpha_n(\overbrace{A, \ldots, A}^{k}, \overbrace{B, \ldots, B}^{n-k})\Bigr|^2 \\
  \leq \Bigl(\sum_k C_n^k\, \alpha_n(\overbrace{|A|, \ldots, |A|}^{k}, \overbrace{|B|, \ldots, |B|}^{n-k})^{\frac12}\, \alpha_n(\overbrace{|A|', \ldots, |A|'}^{k}, |B|', \ldots)^{\frac12}\Bigr)^2
\end{gather*}
\note{les exposants $2$ et $\frac12$ sont surchargés et se lisent avec doute ; les accolades horizontales portent $k$ et $n-k$} d'où, par l'inégalité de Cauchy-Schwartz
\[
  |\alpha_n(A+B)|^2 \leq \Bigl(\sum_k C_n^k\, \alpha_n(|A|, \ldots, |A|, |B|, \ldots, |B|)\Bigr)^{\frac12} \Bigl(\sum_k \ldots\Bigr)^{\frac12}
\]
\note{le second facteur n'est pas écrit en entier : une parenthèse avec $\sum_k$ et une apostrophe ; l'exposant du premier membre est surchargé}

\page{124}
\[
  = \bigl(\alpha_n(|A| + |B|)\bigr)^{\frac12} \bigl(\alpha_n(|A|' + |B|')\bigr)^{\frac12}
\]
\struck{$\det(1 + A + B) \leq \det(1 + |A|)\det(1 + |B|)$}

\emph{Corollaire 3} \note{encadré}
\[
  |\det(1 + A + B)| \leq \tfrac12\bigl(\det(1 + |A| + |B|) + \det(1 + |A|' + |B|')\bigr)
\]
\note{le premier terme de droite se lit « $\det(1 + |A + B|)$ » ou « $\det(1 + |A| + |B|)$ » ; le second est écrit en surcharge} résulte aussitôt de $|\alpha_n(A+B)| \leq \frac12\bigl[\alpha_n(|A| + |B|) + \alpha_n(|A|' + |B|')\bigr]$.

\emph{Corollaire 4} \note{encadré}
\[
  |\det(1 + A + B)| \leq \det(1 + |A|)\det(1 + |B|)
\]
\note{la formule est soulignée d'un double trait}

Il suffit de \uncertain{prouver} \uncertain{ce} \uncertain{résultat} pour $A$ et $B$ hermitiens positifs ; \uncertain{en} \uncertain{effet} \ill{}, d'après le corollaire 3, $|\det(1 + A + B)| \leq \frac12\bigl[\det(1 + |A|)\det(1 + |B|) + \det(1 + |A|')\det(1 + |B|')\bigr]$, ce qui \uncertain{donnera} l'inégalité voulue, en vertu de $\det(1 + |A|) = \det(1 + |A|')$, $\det(1 + |B|) = \det(1 + |B|')$. --- \uncertain{Dans} \uncertain{le} cas \uncertain{hermitien} \uncertain{positif} \note{« positif » interlinéé}, on écrit
\[
  \det(1 + \uncertain{H} + K) = \det(1 + \uncertain{L})\Bigl(1 + \frac{K}{1 + L}\Bigr) = \det(1 + L)\det(1 + HK)
\]
\note{la première lettre se lit $R$ ou $H$, la suite écrit $L$ ; la relation entre les deux n'est pas explicitée} où $H = \dfrac{1}{1+L}$, \uncertain{donc} \uncertain{est} \uncertain{hermitien} \uncertain{positif} \uncertain{de} \uncertain{norme} \ill{}, $\det(1 + HK) \leq \det(1 + K)$, il suffit pour cela de prouver $\alpha_n(HK) \leq \alpha_n(K)$ pour tout $n$, soit $\operatorname{Tr}\Lambda^n HK \leq \|\Lambda^n K\|_1$. \uncertain{Or} \uncertain{on} \uncertain{a}
\[
  \Lambda^n HK = (\Lambda^n H)(\Lambda^n K), \quad \text{d'où} \quad \|\Lambda^n HK\|_1 \leq \|\Lambda^n H\|_\infty\, \|\Lambda^n K\|_1 \quad \text{et} \quad \|\Lambda^n H\|_\infty \leq \|H\|_\infty^{\,n} \leq 1.
\]
\note{l'exposant $n$ de $\Lambda$ est écrit au-dessus de la lettre, comme un accent}

\emph{Corollaire 5} \note{encadré, la formule aussi}
\[
  \det(1 + |A + B|) \leq \det(1 + |A|)\det(1 + |B|)
\]

\page{125}
\note{la partie supérieure de la page reprend la page 124 jusqu'à « il suffit pour cela » ; seul le feuillet collé sur sa partie inférieure est transcrit ici}

Le corollaire \uncertain{contient} \uncertain{le} précédent, car on \uncertain{sait} \uncertain{que} $\det(1 + A + B) \leq \det(1 + |A + B|)$. --- \uncertain{Il} \uncertain{résulte} \uncertain{du} \uncertain{précédent}, car on \uncertain{a} \uncertain{aussi} « \uncertain{en} \uncertain{général} » \note{les guillemets sont de lui} $|A + B| = U(A + B)$, $U$ \uncertain{unitaire}, d'\uncertain{où} \uncertain{par} \uncertain{ce} \uncertain{qui} \uncertain{précède}
\[
  \det(1 + UA + UB) \leq \det(1 + |UA|)\det(1 + |UB|) = \det(1 + |A|)\det(1 + |B|)
\]
(car $|UA| = U|A|U^{-1}$, $|UB| = U|B|U^{-1}$).
\page{126}
\struck{\ill{}} \note{quatre lignes au crayon en tête de page, en partie biffées, dont on lit « Théorème », « le déterminant de Fredholm », « Notation $\mathcal{L}_0$ », « partie du théorème » et « sont connues »}

\section*{Théorème}

1) Pour \struck{\ill{}} $0 < p \leq \infty$, $\mathcal{L}^p(H)$ est un idéal bilatère \note{« bilatère » est interlinéé} dans $\mathcal{L}(H)$, stable pour l'involution $A \to A^*$. De \uncertain{façon} \uncertain{précise} \struck{\ill{}} :
\begin{gather*}
  (1)\qquad \|A\|_p = \|A^*\|_p \qquad \text{pour } A \in \mathcal{L}^p(H),\ B \in \mathcal{L}(H) \\
  (2)\qquad \|BA\|_p \leq \|B\|\,\|A\|_p, \qquad \|AB\|_p \leq \|B\|\,\|A\|_p
\end{gather*}
\note{devant les normes de (1) et (2) un $S_p$ est amorcé puis biffé : il avait d'abord écrit l'énoncé pour $S_p$} Si $\infty \geq p \geq 1$, $A \to \|A\|_p$ est une norme sur $\mathcal{L}^p(H)$ \marginal{\ill{}} i.e. on a
\[
  (3)\qquad \|A + B\|_p \leq \|A\|_p + \|B\|_p, \qquad \|\lambda A\|_p = |\lambda|\,\|A\|_p \qquad (1 \leq p \leq \infty)
\]
\note{le domaine entre parenthèses, en bout de ligne, se lit avec doute} et $\mathcal{L}^p(H)$ est \uncertain{complet} pour cette norme. Si $0 < p \leq \uncertain{1}$, alors l'expression $S_p(A) = \|A\|_p^{\,p}$ sur $\mathcal{L}^p(A)$ \note{ainsi : $\mathcal{L}^p(A)$ pour $\mathcal{L}^p(H)$, ici et deux lignes plus bas} satisfait à l'inégalité du triangle
\[
  (4)\qquad S_p(A + B) \leq S_p(A) + S_p(B) \qquad (0 < p \leq 1)
\]
et muni de la distance, invariante par translation, $S_p(A - B)$, $\mathcal{L}^p(A)$ est un espace vectoriel top. métrisable et complet, non localement convexe si $p < 1$, et $\dim H = +\infty$.

2) Si $A \in \mathcal{L}^p(H)$ ($0 < p < +\infty$), alors la suite des valeurs propres $(\lambda_i)$ \note{interlinéé} de $A$ est de puissance $p$-ième sommable, et on a
\[
  (5)\qquad \sum |\lambda_i|^p \leq S_p(A)
\]
$S_p(A)$ et $\|A\|_p$ sont des fonctions croissantes sur l'ensemble des opérateurs hermitiens positifs, i.e. $H$ et $K$ hermitiens positifs, $H \leq K$, entraîne $S_p(H) \leq S_p(K)$.

3) \note{le chiffre corrige un $2$} Soient $0 < p, q, r \leq +\infty$, avec $\dfrac1p + \dfrac1q \geq \dfrac1r$. Si donc $A \in \mathcal{L}^p(H)$, $B \in \mathcal{L}^q(H)$, alors $AB \in \mathcal{L}^r(H)$, et on a
\[
  (6)\qquad \|AB\|_r \leq \|A\|_p\,\|B\|_q
\]

\marginal{② Dual de $\mathcal{L}^p$ ; réflexivité} \note{en biais dans la marge gauche, à hauteur du point 2)} \marginal{\ill{} \uncertain{généralisation} \ill{} \uncertain{Schatten} \ill{}} \note{seconde note marginale en biais, plus bas, de plusieurs lignes courtes ; une courbe la relie au « (1) ..... » du bas de page} \marginal{(1) .....}

\page{127}
\marginal{\ill{} \uncertain{Démonstration} \uncertain{finie}} \note{dans un cadre, en biais, dans le coin supérieur gauche ; trois lignes dont seules les deux dernières se laissent lire avec doute}

La première partie du théorème est \uncertain{résumée} \uncertain{dans} \uncertain{les} formules (1) à (4), \struck{et} les affirmations relatives \uncertain{aux} \ill{} \uncertain{topologies} \uncertain{naturelles} de \uncertain{ces} \uncertain{espaces} $\mathcal{L}^p(H)$ \struck{\ill{}} \uncertain{d'abord} \uncertain{que} \uncertain{ces} \uncertain{dernières} \uncertain{sont} \uncertain{conséquences} \uncertain{très} \uncertain{faciles} \uncertain{des} \uncertain{formules} \uncertain{(1)} \uncertain{à} \uncertain{(4)} \note{ligne interlinéée en petits caractères}. La formule (1) est évidente, \struck{\ill{}} \note{trois lignes biffées, dont « $p < +\infty$ », « soit », « conséquence » et « Rem. 2) » se lisent} \uncertain{par} \uncertain{ce} \uncertain{que} $|A|$ et $|A^*|$ \uncertain{ont} \struck{\ill{}} \struck{\uncertain{valeurs} \uncertain{propres}} \ill{}. La formule (2) est un cas particulier de la formule (6), mais \struck{\ill{}} \uncertain{demande} \uncertain{une} \uncertain{démonstration} \struck{\ill{}} \struck{\uncertain{directe} \uncertain{simplement} \uncertain{de} \uncertain{façon} \uncertain{directe}} directe, en procédant comme dans [~] \note{les crochets sont vides sur la page}. En effet, si $(k_i)$ \struck{\ill{}} désigne \uncertain{la} suite des valeurs \uncertain{propres} \uncertain{de} $|BA|^{2}$ \note{surchargé ; un « $= (AB)^*(AB)$ » et un « $|BA|$ » biffés, puis « $|A|$ » à la ligne} \marginal{et $(h_i)$ \uncertain{celles} \uncertain{de}} \struck{$|BA|$} rangées par ordre de grandeur décroissante (chacune répétée autant de fois que l'exige sa multiplicité) on a $k_i \leq \|B\|\, h_i$ pour tout $i$ (loc. cit., \uncertain{page} \ill{}) d'où aussitôt les formules (2).

\struck{\ill{}} \note{deux lignes biffées, dont « La \ill{} des formules \ill{} », « (3) peut être regardée », « de la formule », « première \ill{} immédiate »} \uncertain{Les} \uncertain{formules} (3) \uncertain{doivent} \uncertain{être} \uncertain{regardées} \uncertain{comme} \uncertain{conséquence} de la formule (6) (que nous démontrerons de façon indépendante) de la façon suivante. Si $p'$ est tel que $\dfrac1p + \dfrac1{p'} = 1$, et si $B \in \mathcal{L}^{p'}(H)$, \struck{\ill{}} $\|B\|_{p'} \leq 1$, alors on a en vertu de (6) $\|AB\|_1 \leq \|A\|_p$, d'où
\[
  |\operatorname{Tr}(AB)| \leq \|A\|_p, \qquad \text{et} \qquad \operatorname*{Sup}_{\|B\|_{p'} \leq 1} |\operatorname{Tr} AB| \leq \|A\|_p
\]

\page{128}
Cette inégalité peut être remplacée par une égalité, car il suffit \uncertain{en} \uncertain{fait} \uncertain{de} \ill{} \uncertain{isométrie} $U$ \uncertain{telle} \uncertain{que} $UA = |A|$, \struck{\ill{}} et d'\uncertain{autre} \uncertain{part}, \uncertain{comme} \uncertain{il} résulte \uncertain{aussitôt} de la décomposition spectrale de $|A|$ et du th. de Hölder classique sur les espaces $L^p$ et $L^{p'}$, un opérateur hermitien positif $H \in L^{p'}$ tel que $\operatorname{Tr}(H|A|) = \|A\|_p$, soit $\operatorname{Tr}(HUA) = \|A\|_p$. Si on pose $B = \uncertain{HU}$ \note{lu $UA$ ou $HU$ ; le second est ce que demande l'argument}, on aura $\|B\|_{p'} \leq \|U\|\,\|H\|_{p'}$ (formule (2)) donc $\|B\|_{p'} \leq 1$, \uncertain{ce} \uncertain{qui} \uncertain{prouve} \uncertain{notre} \uncertain{assertion}. Par suite on a bien
\[
  (7)\qquad \|A\|_p = \operatorname*{Sup}_{\|B\|_{p'} \leq 1} |\operatorname{Tr} AB|
\]
ce qui prouve que $\|A\|_p$ est bien une semi-norme, comme borne supérieure de semi-normes $A \to |\operatorname{Tr} AB|$.

Une telle démonstration n'est plus praticable pour $p < 1$ ; \struck{et} nous démontrons la formule (4) par voie analytique, grâce au

\emph{Lemme 1} Soit $(\alpha_i)$ une suite \struck{décroissante} de nombres $\geq 0$, \struck{de puissance $p$-ième \ill{}}, \struck{soit} \marginal{et soit} $0 < p < 1$. Posons $f(z) = \prod_{i=1}^{\infty} (1 + \alpha_i z)$. Alors on a
\[
  (8)\qquad \sum \alpha_i^{\,p} = \frac{p \sin \pi p}{\pi} \int_0^\infty \frac{\log f(v)}{v^{1+p}}\, dv
\]
(les deux membres pouvant être infinis simultanément).
\page{129}
\emph{Corollaire} Si $A \in \mathcal{L}^1(H)$, soit $M_A(v) = \operatorname*{Sup}_{|z| = v} \det(1 + zA)$, \struck{si $A$ est hermitien positif, et si $A \in \mathcal{L}^p(H)$, on} \uncertain{et} \uncertain{avec} $0 < p < 1$ ; on a
\[
  (9)\qquad S_p(A) = \frac{p \sin \pi p}{\pi} \int_0^\infty \frac{\log M_{|A|}(v)}{v^{1+p}}\, dv
\]
Il suffit en effet de poser $f(z) = \det(1 + z|A|)$, les $\alpha_i$ étant les valeurs propres de $|A|$ ; on \uncertain{aura} \uncertain{alors} $M(v) = f(v)$.

Pour démontrer (4), on \struck{peut} peut écrire
\[
  S_p(A + B) = S_p(|A + B|) = \frac{p \sin \pi p}{\pi} \int_0^\infty \frac{\log M_{|A+B|}(v)}{v^{1+p}}\, dv
\]
Or $M_{|A+B|}(v) = \det(1 + v|A + B|) \leq \det(1 + v|A|)\det(1 + v|B|)$ (formule \ill{}) \struck{soit $M_A$} soit
\[
  (10)\qquad M_{|A+B|}(v) \leq M_{|A|}(v) \cdot M_{|B|}(v) \qquad \text{d'où}
\]
\[
  S_p(A + B) \leq \frac{p \sin \pi p}{\pi} \int_0^\infty \frac{\log M_{|A|}(v)}{v^{1+p}}\, dv + \frac{p \sin \pi p}{\pi} \int_0^\infty \frac{\log M_{|B|}(v)}{v^{1+p}}\, dv
\]
Appliquant de nouveau (9), on obtient bien la formule voulue (4).

\struck{Pour prouver}

2) Pour prouver la formule (5), je ne connais qu'une méthode analytique. De façon précise, si $A$ est un opérateur cpt \uncertain{quelconque}, \uncertain{soient} la suite de ses val. propres $(\lambda_i)$, \uncertain{posons} \uncertain{et} \uncertain{on} $0 < p < +\infty$, \uncertain{posons}
\[
  s_p(A) = \sum |\lambda_i|^p.
\]
\struck{\ill{}} \note{une ligne biffée, dont « $S_p(A) \leq S_p$ » et « évidemment pour » se lisent} \uncertain{pour} \uncertain{tout} \uncertain{entier} $n > 0$ \note{le début de la ligne est biffé} \quad $s_p(A) = s_{p/n}(A^n)$ \note{$s_{p/n}$ porte l'indice $\frac{p}{n}$ en fraction sous le $s$ ; la ligne est traversée de courbes de renvoi et se termine par un mot perdu}

\page{130}
Supposons d'abord $p < 1$, alors on a, d'après une formule de majoration classique \uncertain{pour} la somme des puissances $p$-ièmes des \uncertain{inverses} \uncertain{des} modules des zéros d'une fonction entière
\[
  f(z) = \det(1 - zA) \qquad (0 < p < 1)
\]
\[
  \struck{(10)}\qquad s_p(A) \leq p^2 \int_0^\infty \frac{\log M_A(r)}{r^{p+1}}\, dr
\]
or $M_A(v) \leq M_{|A|}(v)$ (\ill{}) d'où
\[
  s_p(A) \leq p^2 \int_0^\infty \frac{\log M_{|A|}(v)}{v^{p+1}}\, dv
\]
; utilisant alors la formule (9), on obtient donc la majoration
\[
  (10)\qquad s_p(A) \leq \frac{p\pi}{\sin p\pi}\, S_p(A)
\]
Le facteur $\dfrac{p\pi}{\sin p\pi}$ est $> 1$, mais tend vers $1$ pour $p \to 0$. Soit alors $0 < p < +\infty$ quelconque, et soit $n$ un entier $> 0$ tel que $\dfrac{p}{n} < 1$. On a évidemment $s_p(A) = s_{p/n}(A^n)$, d'où, en vertu de (10) :
\[
  s_p(A) \leq \frac{\frac{p}{n}\pi}{\sin \frac{p}{n}\pi}\, S_{p/n}(A^n)
\]
Or, on a \struck{$S_{p/n}(A^n) \leq S_p(A)$} $\|A^n\|_{p/n} \leq \underbrace{\|A\|_p \cdots \|A\|_p}_{n \text{ facteurs}}$, comme il résulte aussitôt de la formule (6), qui \uncertain{se} démontre plus bas (alors il est immédiat que cette formule se généralise \uncertain{en} \uncertain{prenant} \uncertain{un} \uncertain{nombre} \uncertain{fini} \uncertain{quelconque} \uncertain{d'opérateurs}), d'où
\[
  S_{p/n}(A^n) = \bigl(\|A^n\|_{p/n}\bigr)^{\frac{p}{n}} \leq \|A\|_p^{\,p} = S_p(A)
\]
\page{131}
d'où
\[
  s_p(A) \leq \frac{\pi \frac{p}{n}}{\sin \pi \frac{p}{n}}\, S_p(A).
\]
Faisant tendre $n$ vers l'infini, on trouve l'inégalité (5).

L'assertion que si $0 \leq H \leq K$, \struck{\ill{}} alors on \uncertain{a} \struck{positif} $S_p(H) \leq S_p(K)$ résulte aussitôt du fait bien connu que sous les conditions envisagées, la suite décroissante \note{interlinéé} des valeurs propres de $H$ est majorée par la suite des valeurs propres décroissante de $K$ (\ill{} il résulte de la caractérisation \uncertain{Fischer} classique \note{« classique » interlinéé} \uncertain{extrémale} de la suite des valeurs propres d'un opérateur hermitien complètement continu \uncertain{positif}).

\note{la moitié inférieure de la page, de « 3) » à la fin, est barrée de longs traits obliques ; elle est reprise à la page 132}

3) Pour prouver la formule (6), on \struck{\ill{}} \uncertain{peut} \struck{\ill{}} \struck{écrire, pour un entier $n > 0$ arbitraire} $S_{r/2}(L)$ où $L = (AB)(AB)^*$ \note{encadré}
\[
  S_r(AB) = S_{r/2}\bigl(AB(AB)^*\bigr) = \struck{S_{r/2m}\bigl((AB(AB)^*)^m\bigr)}
\]
Si $m$ est un entier $> 0$, on \uncertain{aura} \uncertain{encore}
\[
  S_r(AB) = S_{r/2}(L) = S_{r/2m}(L^m), \quad \text{d'où, en vertu de la formule (9),} \qquad k(\alpha) = \frac{\pi\alpha}{\sin \pi\alpha}
\]
\note{$k(\alpha)$ est défini à droite, dans un renvoi ; suivent deux lignes encadrées et biffées, dont on lit « $r' = \frac{r}{2m}$ » et « $S_r(AB) = \frac{\ill{}}{k(r/2)} \int_0^\infty \frac{\log M_{L^m}(v)}{v^{1+r/2m}}\, dv$ »} Posons $r' = \dfrac{r}{2m}$, $p' = \dfrac{p}{2m}$, $q' = \dfrac{q}{2m}$, on \uncertain{aura} évidemment encore $\dfrac{1}{p'} + \dfrac{1}{q'} \geq \dfrac{1}{r'}$. La formule (9) donne alors
\[
  S_r(AB) = S_{r'}(L^m) = \frac{r' \sin \pi r'}{\pi} \int_0^\infty \frac{\log M_{L^m}(v)}{v^{1+r'}}\, dv
\]

\page{132}
3) La preuve de la formule (6) est la plus difficile. Nous commencerons par donner une preuve directe et très simple pour le cas où $p \geq 1$, $q \geq 1$, $r = 1$ \note{interlinéé au-dessus d'un « $p, q, r \geq$ \ill{} » biffé, avec un astérisque} (preuve qui, d'ailleurs, \uncertain{a} \uncertain{l'avantage} \uncertain{de} \uncertain{s'appliquer} \uncertain{sans} modification au cas d'une \struck{D} \uncertain{trace} \uncertain{au} \uncertain{sens} \uncertain{de} \uncertain{von} \uncertain{Neumann} \uncertain{définie} sur une algèbre auto-adjointe $\mathfrak{A}$ \struck{\ill{}} \marginal{(1)} (\struck{les éléments} les « \uncertain{opérateurs} \struck{\ill{}} de puissance $p$-ième \uncertain{sommable} » de $\mathfrak{A}$ se définissent comme \uncertain{ceux} \uncertain{des} $A \in \mathfrak{A}$ tels que $\struck{\operatorname{Tr}}(|A|^p) < +\infty$.) \note{une accolade relie ces lignes à une note marginale en biais, illisible hormis « Schatten » et un point d'interrogation redoublé} (\ill{} \uncertain{obtiendrons} d'ailleurs \uncertain{le} \uncertain{résultat} comme conséquence \uncertain{d'une} proposition plus générale).

Nous commencerons par prouver que si $A_1, \ldots, A_{2^n}$ sont $2^n$ opérateurs cpt \uncertain{continus} ($n$ entier $\geq 1$) alors on a ${}^{(1)}$
\[
  (11)\qquad \|A_1 A_2 \cdots A_{2^n}\|_1 \leq \|A_1\|_{2^n}\, \|A_2\|_{2^n} \cdots \|A_{2^n}\|_{2^n}
\]
\note{les $2^n$ sont écrits « $2n$ » sur la ligne, ici et dans la suite ; la récurrence sur $n$ avec $2^{n-1}$ facteurs fixe la lecture} Nous le prouvons par récurrence sur $n$ ; c'est vrai si $n = 1$, \struck{(formule \ill{} bien), \ill{}} \struck{$\|A_1 A_2\|_1 \leq \|A_1\|_2 \|A_2\|_2$} \uncertain{supposons-le} \uncertain{démontré} pour $2^{n-1}$ facteurs, et démontrons alors (11). Écrivons $A_1 \cdots A_{2^n} = (A_1 A_2) \cdots (A_{2^n - 1} A_{2^n})$ ; le premier membre \uncertain{est}, en vertu de l'hypothèse de récurrence, majoré par $\|A_1 A_2\|_{2^{n-1}} \cdots \|A_{2^n - 1} A_{2^n}\|_{2^{n-1}}$

\note{au bas de la page, sous un trait, la note de renvoi :} ${}^{(1)}$ Le lecteur notera que ce résultat suffirait pour démontrer \struck{directement} directement \note{réécrit au-dessus du mot biffé}, \uncertain{comme} \uncertain{nous} \uncertain{le} \uncertain{verrons} plus bas, la formule (6) dans le cas général.
\page{133}
\uncertain{On} \uncertain{voit} \uncertain{qu'on} est ramené à démontrer \uncertain{la} formule \struck{telle} \uncertain{pour} $\|A_1 B_1\|_{2^{n-1}} \leq \|A_1\|_{2^n} \|B_1\|_{2^n}$, soit, \uncertain{plus} \uncertain{simplement}, $\|AB\|_{2^{n-1}} \leq \|A\|_{2^n}\, \|B\|_{2^n}$. \uncertain{On} \uncertain{a} \ill{}
\[
  \|AB\|_{2^{n-1}} = \Bigl(\bigl\|\bigl(AB(AB)^*\bigr)^{\frac{2^{n-1}}{2}}\bigr\|_1\Bigr)^{\frac{1}{2^{n-1}}} = \bigl(\|C\|_1\bigr)^{\frac{1}{2^{n-1}}}
\]
où
\[
  C = (ABB^*A^*)^{2^{n-2}} = A(BB^*)(A^*A)\cdots(BB^*)A^*
\]
($2^{n-1} + 1$ facteurs \uncertain{dans} \uncertain{le} dernier membre) \uncertain{on} \uncertain{a}
\[
  \|C\|_1 = \operatorname{Tr} C = \operatorname{Tr} C' \quad \text{où } C' = A^*A \cdot BB^* \cdot A^*A \cdots BB^*
\]
\note{« $= \operatorname{Tr} C'$ » est interlinéé au-dessus d'un membre biffé ; un « $= \operatorname{Tr}$ » biffé en fin de ligne} \struck{\ill{} dernier membre figurant $2^{n-1}$} facteurs ($2^{n-1}$ facteurs \uncertain{dans} \uncertain{le} second membre). Or $|\operatorname{Tr} C'| \leq \|C'\|_1$, et d'après l'hypothèse de récurrence, on a \uncertain{de} \uncertain{nouveau}
\[
  \|C'\|_1 \leq \|A^*A\|_{2^{n-1}}\, \|BB^*\|_{2^{n-1}} \cdots \|BB^*\|_{2^{n-1}} = \bigl(\|A^*A\|_{2^{n-1}}\bigr)^{2^{n-2}} \bigl(\|B^*B\|_{2^{n-1}}\bigr)^{2^{n-2}}
\]
Or $\|A^*A\|_{2^{n-1}} = \bigl(\|A\|_{2^n}\bigr)^2$, $\|B^*B\|_{2^{n-1}} = \bigl(\|B\|_{2^n}\bigr)^2$, d'où $\|C'\|_1 \leq \bigl(\|A\|_{2^n} \|B\|_{2^n}\bigr)^{2^{n-1}}$. On obtient donc bien
\[
  \|AB\|_{2^{n-1}} = \bigl(\|C\|_1\bigr)^{\frac{1}{2^{n-1}}} \leq \bigl(\|C'\|_1\bigr)^{\frac{1}{2^{n-1}}} \leq \|A\|_{2^n}\, \|B\|_{2^n},
\]
ce qui prouve la formule (11).

\page{134}
\struck{4)} Soient \struck{$p \geq 1$, $q \geq 1$} $\dfrac1p + \dfrac1q = 1$, $A \in \mathcal{L}^p$, $B \in \mathcal{L}^q$, alors $AB \in \mathcal{L}^1$, et $\|AB\|_1 \leq \|A\|_p\, \|B\|_q$. \uncertain{Preuve} \uncertain{On} \uncertain{peut} \uncertain{supposer} $A, B \in \mathcal{L}^1$ \note{ligne interlinéée, « Preuve » souligné}. On a $A = UH$, $B = KV$ ($U$, $V$ partiellement isométriques). Posons \struck{\ill{}} $n$, soient $\alpha, \beta$ deux entiers \ill{} \uncertain{tels} \uncertain{que} $\alpha + \beta = 2^n$, posons $h = H^{\frac1\alpha}$, $k = K^{\frac1\beta}$, \ldots $AB = U h^\alpha k^\beta V$ d'où $\|AB\|_1 \leq \|h^\alpha k^\beta\|_1 \leq$ \uncertain{or} (Lemme 1) \note{ainsi ; c'est la formule (11) qui est invoquée}
\[
  \|h^\alpha k^\beta\|_1 \leq \underbrace{\|h\|_{2^n} \cdots \|h\|_{2^n}}_{\alpha} \underbrace{\|k\|_{2^n} \cdots \|k\|_{2^n}}_{\beta} = \bigl(\operatorname{Tr} h^{2^n}\bigr)^{\frac{\alpha}{2^n}} \bigl(\operatorname{Tr}(k^{2^n})\bigr)^{\frac{\beta}{2^n}}
\]
d'où
\[
  \|AB\|_1 \leq \bigl(\operatorname{Tr} H^{\lambda}\bigr)^{\frac1\lambda} \bigl(\operatorname{Tr} K^{\mu}\bigr)^{\frac1\mu} \qquad \text{où } \lambda = \frac{2^n}{\alpha},\ \mu = \frac{2^n}{\beta}
\]
\note{l'exposant de $H$ est surchargé, $\frac{2^n}{\alpha}$ corrigé en $\lambda$} \struck{Faisons alors} \uncertain{on} \uncertain{voit} Cette inégalité est vraie pour ($\lambda, \mu > 1$) \uncertain{rationnels} \uncertain{tels} \uncertain{qu'on} \uncertain{puisse} \uncertain{les} mettre sous \uncertain{la} forme $\dfrac{2^n}{\alpha} \cdot \dfrac{2^n}{\beta}$ avec $\alpha + \beta = 2^n$, i.e. pour \uncertain{tous} $\dfrac1\lambda + \dfrac1\mu = 1$, \struck{\ill{}} et $\dfrac1\lambda$ et $\dfrac1\mu$ \uncertain{sont} \uncertain{des} multiples de $\dfrac{1}{2^n}$ ; par un passage à la limite \uncertain{évident}, on \uncertain{conclut}
\[
  \|AB\|_1 \leq \bigl(\operatorname{Tr} H^p\bigr)^{\frac1p} \bigl(\operatorname{Tr} H^{p'}\bigr)^{\frac{1}{p'}}
\]
\note{ainsi, $H^{p'}$ pour $K^{p'}$} \marginal{\ill{} $\|A_1 \cdots A_n\|_1 \leq \|A_1\|_{p_1} \cdots$ \uncertain{si} $\dfrac{1}{p_1} + \dfrac{1}{p_2} + \cdots + \dfrac{1}{p_n} = 1$} \note{interlinéé en petits caractères, avec une flèche partant de la marge gauche ; la généralisation à $n$ facteurs}

3) \uncertain{On} \uncertain{en} \uncertain{conclut} \uncertain{que} \uncertain{l'on} \uncertain{a} ($p \geq 1$, $\dfrac1p + \dfrac1{p'} = 1$)
\[
  \|A\|_p = \operatorname*{Sup}_{\|B\|_{p'} \leq 1} |\operatorname{Tr} AB| \qquad (\text{d'où} \ \ill{})
\]
\[
  \|A_1 + A_2\|_p \leq \|A_1\|_p + \|A_2\|_p \qquad ) \ \uncertain{de} \ \uncertain{même}
\]
\struck{$AB$} si $\dfrac1p + \dfrac1q \geq \dfrac1r$ \quad ($p \geq 1$, $q \geq 1$, $r \geq 1$)
\page{135}
Le résultat précédent nous permet de démontrer le résultat auxiliaire suivant, qui nous permettra de démontrer la formule (6) en toute généralité :

\emph{Lemme} Soient $A$ et $B$ deux opérateurs $\in \mathcal{L}^{2m}(H)$, \marginal{posons} $H = A^*A$, $K = B^*B$, $L = (AB)^*(AB)$. Soient \struck{et soit} $m$ et $n$ des entiers $\geq 1$. Alors \struck{$L^n$, $H^n$} $L^m$, $H^m$ et $K^m$ \note{les exposants sont surchargés, $n$ corrigé en $m$, et « $L^m$, $H^m$ » repris dans la marge droite} \uncertain{sont} \uncertain{opérateurs} à trace, et on a
\[
  (12)\qquad \alpha_n(L^m) \leq \alpha_n(H^m)\, \alpha_n(K^m)
\]
Le fait que $L^m$, $H^m$, $K^m$ soient opérateurs à trace résulte \uncertain{immédiatement} \ill{} \uncertain{de} \uncertain{ce} \uncertain{qui} \uncertain{précède} ; \uncertain{si} \uncertain{l'on} \uncertain{vérifie} (12) \uncertain{s'écrit} \uncertain{alors}, \struck{et} en posant $A' = \Lambda^n A$, $B' = \Lambda^n B$ \note{deux lignes serrées, en partie biffées ; le passage à $\Lambda^n$ est ce qui ramène (12) au cas $n = 1$} 
\begin{gather*}
  \struck{H = A^*A,\ K = B^*B,\ L = (AB)^*(AB)} \\
  \struck{\operatorname{Tr} L^m \leq \operatorname{Tr} H^m\, \operatorname{Tr} K^m} \\
  \struck{(\|L\|_m)^m \leq (\|H\|_m)^m (\|K\|_m)^m} \\
  \struck{\|AB\|_m \leq \|A\|_m \|B\|_m.\ \text{Or, on a}\ \ill{}} \\
  \struck{\ill{}\ \|AB\|_m \leq \|A\|_{\frac m2} \|B\|_{\frac m2}}
\end{gather*}

\note{cinq lignes biffées, la deuxième et la troisième encadrées d'une accolade} ce qui nous ramène au cas où $n = 1$. On \uncertain{doit} alors démontrer
\[
  \operatorname{Tr}\bigl[(ABB^*A^*) \cdots (ABB^*A^*)\bigr] \leq \bigl(\|A^*A\|_m\bigr)^m \bigl(\|B^*B\|_m\bigr)^m
\]
\note{le premier membre porte un indice $1$ ou $m$ après le crochet, surchargé} Or, le premier membre \struck{\ill{} la même chose \ill{}} \struck{d'un produit de $4m$ facteurs est \ill{}} \uncertain{est} \uncertain{la} \uncertain{trace} du produit de $2m$ facteurs \struck{majoré par} $(A^*A)(BB^*)\cdots(BB^*)$, donc \uncertain{est} \uncertain{majoré} \uncertain{par} \uncertain{la} \uncertain{norme} \uncertain{de} \uncertain{ce} produit, donc \uncertain{en} \uncertain{vertu} \uncertain{de} \uncertain{(11)} :
\[
  \bigl(\|A^*A\|_{2m}\bigr)^m \bigl(\|BB^*\|_{2m}\bigr)^m
\]
d'après la formule $\|A_1 \cdots A_{2m}\|_1 \leq \|A_1\|_{2m} \cdots \|A_{2m}\|_{2m}$. Or, on a $\|A^*A\|_{2m} \leq \|A^*A\|_m$, et $\|BB^*\|_{2m} = \|B^*B\|_{2m} \leq \|B^*B\|_m$, d'où résulte le lemme.

\page{136}
Soient maintenant $A \in \mathcal{L}^p$, $A \in \mathcal{L}^q$ \note{ainsi, $A$ pour $B$}, \marginal{$0 < p, q, r \leq \infty$} \note{interlinéé au-dessus de la première ligne} $\dfrac1p + \dfrac1q \geq \dfrac1r$, démontrons la formule (6). Nous pouvons supposer \uncertain{que} $p$ et $q$ \uncertain{sont} \uncertain{finis} \struck{\ill{} si l'un \uncertain{d'eux} était infini, \ill{} les formules (2)} \struck{\ill{} deux infinis (\ill{})} \note{deux lignes biffées, la première interlinéée} \uncertain{et} \uncertain{...} démontrons le résultat. \struck{$\|AB\|_r \leq \|AB\|_{\ill{}} = \|A\|_{\ill{}} \|B\|_{\ill{}}$ \ill{}} De plus, \uncertain{on} \uncertain{peut} \uncertain{supposer} \uncertain{évidemment} \struck{\ill{}} \uncertain{que} l'on a \uncertain{exactement} $\dfrac1p + \dfrac1q = \dfrac1r$ (puisque $\|A\|_r$ est fonction décroissante de $r$). Soit $m$ un entier tel que $p \leq 2m$, $q \leq 2m$. \struck{Alors $A$ \ill{} $L$ \ill{}} \struck{On a} \struck{$S_r(AB) = S_{r/2}(AB(AB)^*) = S_{r/2m}(L)$} $L = (AB)^*(AB)$, $H = A^*A$, $K = B^*B$, \uncertain{on} \uncertain{aura}
\[
  S_r(AB) = S_{r/2}(L) = S_{r/2m}(L^m) \struck{\ill{}}
\]
d'où, en vertu de la formule (9), on \uncertain{trouve}, \uncertain{en} \uncertain{posant} $\lambda(x) = \dfrac{\pi}{\sin \pi x}$
\[
  (11)\qquad S_r(AB) = S_{r/2m}(L^m) = \frac{r/2m}{\lambda(r/2m)} \int_0^\infty \frac{\log M_{L^m}(\rho)}{\rho^{1 + r/2m}}\, d\rho
\]
\note{le numéro (11) est repris ici pour une autre formule que celle de la page 132} Or, \struck{$M$} $M_{L^m}(\rho) = \sum \alpha_n(L^m)\, \rho^n$, et \struck{$M_{L^m}(\rho) = \sum_n \alpha_n(L^m)\rho^n$} $\alpha_n(L^m) \leq \alpha_n(H^m)\, \alpha_n(K^m)$ en vertu du lemme 2 \struck{$\alpha_n(L^m) \leq \alpha_n(H^m) \ill{}$}, d'où
\[
  M_{L^m}(\rho) \leq \sum \alpha_n(H^m)\, \alpha_n(K^m)\, \rho^n
\]
Posons $\Bigl[\, p' = \dfrac1p,\ q' = \dfrac1q,\ r' = \dfrac1r \,\Bigr]$, on a : \struck{$p' + q' = r'$} $p'r + q'r = 1$, d'où
\[
  M_{L^m}(\rho) \leq \sum \alpha_n(H^m)\, \alpha_n(K^m)\, \rho^{np'r + nq'r} = \sum_n \bigl(\alpha_n(H^m)\rho^{np'r}\bigr)\bigl(\alpha_n(K^m)\rho^{nq'r}\bigr)
\]
\marginal{\ill{}} \note{trois lignes en biais dans la marge gauche, à hauteur des deux dernières formules, illisibles}
\page{137}
\[
  \leq \Bigl(\sum_n \alpha_n(H^m)\, \rho^{np'r}\Bigr)\Bigl(\sum_n \alpha_n(K^m)\, \rho^{nq'r}\Bigr) \struck{\leq} = M_{H^m}(\rho^{p'r})\, M_{H}(\rho^{q'r})
\]
\note{ainsi, $M_H$ pour $M_{K^m}$} d'où, compte tenu de (12),
\[
  (13)\qquad \frac{\lambda(r/2m)}{r/2m}\, S_r(AB) \leq \int_0^\infty \frac{\log M_{H^m}(\rho^{p'r})}{\rho^{1 + r/2m}}\, d\rho + \int_0^\infty \frac{\log M_{K^m}(\rho^{q'r})}{\rho^{1 + r/2m}}\, d\rho
\]
Calculons \struck{par exemple} la première intégrale du second membre, par le changement de variable $\sigma = \rho^{p'r}$ d'où $\rho = \sigma^{pr'}$, $\dfrac{d\rho}{\rho} = pr'\,\dfrac{d\sigma}{\sigma}$. il vient
\[
  \int_0^\infty \frac{\log M_{H^m}(\rho^{p'r})}{\rho^{r/2m}}\, \frac{d\rho}{\rho} = pr' \int_0^\infty \frac{\log M_{H^m}(\sigma)}{\sigma^{p/2m}}\, \frac{d\sigma}{\sigma}
\]
Le deuxième membre est égal (formule (9)) à
\[
  pr'\, \frac{\lambda(p/2m)}{p/2m}\, S_{p/2m}(H^m) = \frac{\lambda(p/2m)}{r/2m}\, S_p(A).
\]
\note{le numérateur du dernier membre est surchargé} \struck{De} \uncertain{De} \uncertain{même} \uncertain{la} deuxième intégrale dans (13) est égale à $\dfrac{\lambda(q/2m)}{r/2m}\, S_q(B)$, et (13) s'écrit \struck{\ill{}} :
\[
  \lambda(r/2m)\, S_r(AB) \leq \lambda(p/2m)\, S_p(A) + \lambda(q/2m)\, S_q(B)
\]
\uncertain{ou} \uncertain{encore}
\[
  S_r(AB) \leq \frac{\sin \pi \frac{r}{2m}}{\sin \pi \frac{p}{2m}}\, S_p(A) + \frac{\sin \pi \frac{r}{2m}}{\sin \pi \frac{q}{2m}}\, S_q(B)
\]
\note{une ligne interlinéée en très petits caractères, dont on lit « cette relation \ill{} $S_r(AB) < +\infty$ » et « $S_p(A)$ et $S_q(B)$ » : elle observe que l'inégalité donne déjà la finitude de $S_r(AB)$} Faisant tendre $m$ vers $+\infty$, on obtient \uncertain{la} formule plus simple
\[
  (14)\qquad S_r(AB) \leq \frac{r}{p}\, S_p(A) + \frac{r}{q}\, S_q(B)
\]
\note{encadrée ; le $p$ du premier dénominateur est surchargé} C'est une inégalité remarquable, mais n'est pas celle que nous cherchons. Mais \struck{cette relation} \struck{\ill{}}

\page{138}
\struck{\uncertain{à} \uncertain{part} d'homogénéité évidente, qui permet} Remplaçons dans cette formule $A$ par $\lambda A$, $B$ par $\mu B$ ($\lambda$ et $\mu$ nombres $\geq 0$) on obtient
\[
  \struck{\ill{}}\qquad (\lambda\mu)^r\, S_r(AB) \leq \frac{\lambda^p}{p}\, S_p(A) + \frac{\mu^q}{q}\, S_q(B)
\]
\note{le premier membre est réécrit sous une première version biffée} \struck{Exprimons que cette inégalité doit être vérifiée pour \ill{} $\lambda, \mu \geq 0$.} Pour simplifier, posons
\[
  a = S_p(A), \quad b = S_q(B), \quad c = S_r(C), \quad \ill{}
\]
\note{ainsi, $S_r(C)$ pour $S_r(AB)$}
\[
  c \leq \operatorname*{Inf}_{\substack{\lambda \geq 0 \\ \mu \geq 0}} \frac{a\,\dfrac{\lambda^p}{p} + b\,\dfrac{\mu^q}{q}}{\struck{\ill{}}\,(\lambda\mu)^r}
\]
\struck{Posons $x = \lambda^p$, $y = \mu^q$} Posons $x = \lambda^p$, $y = \mu^q$, $\alpha = \dfrac{r}{p}$, $\beta = \dfrac{r}{q}$, d'où \marginal{$\dfrac rp + \dfrac rq = 1$} \note{souligné, dans la marge droite} \uncertain{on} \uncertain{a}
\[
  c \leq \operatorname*{Inf}_{\substack{x \geq 0 \\ y \geq 0}} \frac{\alpha a\, x + \beta b\, y}{x^\alpha y^\beta} = \operatorname*{Inf}_{\substack{x \geq 0 \\ y \geq 0 \\ x^\alpha y^\beta = 1}} \alpha a\, x + \beta b\, y
\]
\note{le numérateur porte des exposants biffés : il avait d'abord écrit $\alpha a\, x^{\ill{}} + \beta b\, y^{\ill{}}$} \marginal{$c \leq a^{\uncertain{r/p}} b^{\uncertain{r/q}}$ \quad $c \leq a^\alpha b^\beta$} \note{en biais dans la marge gauche, sous un trait ; c'est le résultat visé} On vérifie aussitôt que le \struck{minimum} \uncertain{borne} \uncertain{inférieure} du second membre est atteinte (car l'on a \uncertain{points} $(x, y)$ tels que $x \geq 0$, $y \geq 0$, $\alpha a x + \beta b y \leq M$ \ill{} \uncertain{est} \uncertain{compact}) \note{la parenthèse est écrite sur deux lignes serrées, la seconde interlinéée} elle est donc atteinte en un point $(x, y)$ \uncertain{tel} \uncertain{que} \uncertain{les} \uncertain{équations}, \struck{$\alpha a x + \beta b y = 0$} $d(\alpha a x + \beta b y) = 0$ \uncertain{aient} \uncertain{une} \uncertain{solution} \uncertain{non} \uncertain{nulle} \uncertain{en} $d(x^\alpha y^\beta) = 0$ \struck{soient compatibles} \uncertain{en} $dx$, $dy$ ; \uncertain{or} $d(x^\alpha y^\beta) = \alpha x^{\alpha-1} y^\beta\, dx + \beta y^{\beta-1} x^\alpha\, dy$ ; i.e. \struck{tel que $d(\ill{})$} \uncertain{la} \uncertain{condition} d'extremum est donc
\[
  \frac{\alpha a}{\alpha x^{\alpha-1} y} = \frac{\beta b}{\beta x^\alpha y^{\beta-1}}
\]
\note{ainsi ; le $y$ du premier dénominateur est sans exposant} soit $ax = by$ \struck{\ill{}} ; \struck{d'où $ax = by = \dfrac{\alpha a x + \beta b y}{\alpha + \beta}$} ; \struck{$\alpha a x + b y = \dfrac{\alpha a x}{\ill{}}$ \ill{}} \uncertain{comme} $\alpha + \beta = 1$, le minimum cherché \uncertain{est} \uncertain{donc} \struck{$\alpha x + b$} \uncertain{égal} \uncertain{à} $ax = by$ ; \uncertain{compte} \uncertain{tenu} \uncertain{de} \uncertain{la} \uncertain{condition}

\page{139}
\uncertain{d'où} $x^\alpha y^\beta = 1$ \uncertain{on} \uncertain{trouve} $x = \Bigl(\dfrac ba\Bigr)^\beta$, d'où \struck{\ill{}} \uncertain{le} \uncertain{minimum} \uncertain{cherché} $ax = b^\beta a^{1-\beta} = a^\alpha b^\beta$. \note{le $b^\beta$ est corrigé au-dessus d'un facteur biffé} On trouve donc $c \leq a^\alpha b^\beta$, \uncertain{ce} \uncertain{qui} \uncertain{achève} la démonstration.

\end{document}
